% ===========================================================================
\chapter{Words on the map}\label{ch:labels}

A map without names is a puzzle. This chapter puts the names of places,
areas and roads on the map, which turns out to be one of the hard problems
of map making: there are more names than there is room for, and a label
that overlaps another makes both unreadable.

\section{Placing labels greedily}

Professional map makers place labels with care, trying many positions for
each and weighing them all against each other. We'll use the simplest
method that works well, which is also what many map apps do when they have
to be quick:

\begin{enumerate}[nosep]
  \item Put the labels in order of importance.
  \item Take them one at a time. Work out the rectangle the label would
    cover. If it overlaps a label already placed, skip it; otherwise place
    it, and its rectangle is taken.
\end{enumerate}

The important names always make it onto the map, and the less important
ones fill whatever room is left. Zoom in, and there is more room, so more
names appear. That is exactly what happens on a phone as you pinch to
zoom.

Our order is: place names, then area names, then road names, with main
roads before lanes before streets before tracks and paths. For a road, we
try three spots: halfway along it, then a quarter and three quarters of
the way.

We need a rectangle that can tell whether it overlaps another. Two
rectangles overlap when they overlap both left to right and bottom to
top, and two ranges overlap when each one starts before the other ends.
And a label: some text, where it goes, and how it looks.

\begin{addbefore}{map_view.nex}{before the class Map_View}{class Map_View}
-- A rectangle on the map, for keeping labels apart.
class Box
create
  make(l, b, r, t: Real) do
    left := l
    bottom := b
    right := r
    top := t
  end
feature
  once left: Real
  once bottom: Real
  once right: Real
  once top: Real

  overlaps(other: Box): Boolean do
    result := left < other.right and other.left < right
      and bottom < other.top and other.bottom < top
  end
end

-- A piece of text to put on the map.
class Label
create
  make(p: Point, s: String, f: Font, c: Color, in_middle: Boolean) do
    at := p
    text := s
    font := f
    colour := c
    centred := in_middle
  end
feature
  once at: Point
  once text: String
  once font: Font
  once colour: Color
  -- Whether the text is centred on `at`, or starts there.
  once centred: Boolean
end
\end{addbefore}

The view keeps the labels that have been placed and the room they take,
and the three fonts it uses:

\begin{addbefore}{map_view.nex}{after the field metres_per_pixel}{  -- Pixels at the top and bottom of the canvas that something else covers}
  -- The labels that have found room on the map, and the room they take.
  labels: Array[Label]
  taken: Array[Box]
  PLACE_FONT = Font.SANS.bold().with_size(12)
  AREA_FONT = Font.SERIF.italic().with_size(13)
  ROAD_FONT = Font.SANS.with_size(10)
\end{addbefore}

\begin{tightaddbefore}{map_view.nex}{to make, before show_all()}{    show_all()}
    labels := []
    taken := []
\end{tightaddbefore}

\section{Placing the labels}

Places are labelled just to the right of their dot, areas in the middle of
the area (the average of its corners), and roads at the three spots along
them. Each one goes through \code{add_label}, which decides whether it
fits.

\begin{addbefore}{map_view.nex}{before draw_scale_bar}{  -- A bar of a round length (100 m, 200 m, 500 m, 1 km, ...) in the}
  -- Labels must not overlap, or none of them can be read. The map places
  -- them greedily, most important first: each label goes in only if there
  -- is still room for it, and takes that room from the labels after it.
  -- Place names come first, then area names, then road names, main roads
  -- before lanes before streets.

  -- Each place's name, just to the right of its dot.
  place_place_labels() do
    across island.places as p do
      let at: Point := island.point(p.node_id) + create Point.make(metres(8.0), 0.0)
      add_label(create Label.make(at, p.name, PLACE_FONT, Palette.LABEL, false))
    end
  end

  -- Area names in the middle of each area (except the island itself,
  -- whose name is the map's title).
  place_area_labels() do
    across island.areas as a do
      if a.kind /= "land" then
        add_label(create Label.make(average(island.points_of(a.node_ids)), a.name,
          AREA_FONT, when a.kind = "water" then Palette.WATER_LABEL else Palette.LABEL end,
          true))
      end
    end
  end

  -- Each road's name, halfway along it if there is room there, or else a
  -- quarter or three quarters of the way along.
  place_road_labels() do
    across Road_Kind.values as kind do
      across island.roads as r do
        if r.kind = kind then
          let pts: Array[Point] := island.points_of(r.node_ids)
          let length: Real := line_length(pts)
          let placed: Boolean := false
          across [0.5, 0.25, 0.75] as fraction do
            if not placed then
              placed := add_label(create Label.make(point_along(pts, length * fraction),
                r.name, ROAD_FONT, Palette.ROAD_LABEL, true))
            end
          end
        end
      end
    end
  end
\end{addbefore}

\code{add_label} needs to know how much room a label takes. The canvas can
measure text exactly, but only privately, while drawing it; we estimate
instead. In these fonts, a character is on average a little more than half
as wide as the font's size, so a label of $n$ characters in a 10-point font
is about $0.58 \times 10 \times n$ pixels wide, and a few pixels of padding
around it keeps neighbours from touching.

A label whose rectangle is entirely off the canvas is skipped too: it
can't be seen, and it would only take room that doesn't exist.

\begin{continued}{map_view.nex}
  -- Add `label` if there is room for it, and answer whether there was.
  add_label(label: Label): Boolean do
    let w: Real := metres(label.text.length.to_real * label.font.size.to_real * 0.58 + 6.0)
    let h: Real := metres(label.font.size.to_real + 4.0)
    let left: Real := when label.centred then label.at.x - w / 2.0 else label.at.x end
    let room: Box := create Box.make(left, label.at.y - h / 2.0, left + w, label.at.y + h / 2.0)
    -- A label off the canvas takes no room, and isn't drawn.
    result := room.right > canvas.x_min and room.left < canvas.x_max
      and room.top > canvas.y_min and room.bottom < canvas.y_max
    across taken as other do
      if room.overlaps(other) then
        result := false
      end
    end
    if result then
      taken.add(room)
      labels.add(label)
    end
  end
\end{continued}

\section{Halos}

A label drawn straight onto the map is hard to read where it crosses a
road or the edge of a wood. Map makers draw a \emph{halo} around the
letters: a thin outline in the background colour, which separates them
from whatever is underneath. We make one by drawing the text in white
eight times, nudged one pixel in each direction, and then once in its own
colour on top.

That is nine pieces of text for every label, and the pen colour changes
twice for each. We saw in the last chapter that changing the pen is slow,
so we draw all the halos first, in one colour, then all the text, and
change the font or colour only when it differs from the label before.

\begin{continued}{map_view.nex}
  -- Draw the labels that found room: first a white outline around every
  -- one, so that they stay readable over roads and areas, then the text
  -- itself.
  draw_labels() do
    canvas.set_pen_color(Palette.HALO)
    across labels as label do
      use_font(label.font)
      across halo_offsets() as o do
        label_text(label, label.at + o)
      end
    end
    across labels as label do
      use_font(label.font)
      if label.colour /= canvas.pen_color then
        canvas.set_pen_color(label.colour)
      end
      label_text(label, label.at)
    end
  end

  -- Labels come in only a few fonts and colours, and the canvas is slow
  -- to change them (see draw_roads), so they are changed only when they
  -- differ.
  use_font(f: Font) do
    if f /= canvas.font then
      canvas.set_font(f)
    end
  end

  label_text(label: Label, p: Point) do
    if label.centred then
      canvas.text(p.x, p.y, label.text)
    else
      canvas.text_left(p.x, p.y, label.text)
    end
  end

  -- One pixel's worth of nudge in eight directions.
  halo_offsets(): Array[Point] do
    let d: Real := metres(1.2)
    across [-1.0, 0.0, 1.0] as i do
      across [-1.0, 0.0, 1.0] as j do
        if i /= 0.0 or j /= 0.0 then
          result.add(create Point.make(i * d, j * d))
        end
      end
    end
  end
\end{continued}

Last, \code{draw_map} places and draws the labels after the places' dots,
starting each time with no labels and no room taken:

\begin{tightaddbefore}{map_view.nex}{to draw_map, before draw_scale_bar()}{    draw_scale_bar()}
    taken := []
    labels := []
    place_place_labels()
    place_area_labels()
    place_road_labels()
    draw_labels()
\end{tightaddbefore}

\section{The finished map}

\begin{tryit}{try_labels.nex}
intern Map_View

let island: Island_Map := create Island_Map.load(user_path("kestrel_island.map"))
let view: Map_View := create Map_View.make(island, 1000, 700)
view.draw_map()
view.save("island.png")
say(view.labels.length.to_string + " labels found room")

-- Zoom in on the town: 1.4 metres per pixel, centred on Market Square.
if ?island.find_place("market") as market then
  view.look_at(island.point(market.node_id), 1.4)
  view.draw_map()
  view.save("town.png")
  say(view.labels.length.to_string + " labels found room in the town")
end
\end{tryit}
\tryoutput{try_labels}
\mappicture{island}

At this scale there is room for every place name, every area but the town
(whose name would cover Market Square's), and most of the roads. In the
town, the streets are too close together for all their names. Zoomed in,
all of them fit, and the one-way arrows appear:

\mappicture{town}

The town has four one-way streets: Market Street runs east and Church
Street west, so that between them traffic can go both ways, and the same
for Net Lane (north) and Chapel Street (south). They will make the
driving directions in town more interesting than the walking ones.
